% =====================================================================
%  COVER LETTER -- Radiology: Artificial Intelligence, Original Research
%
%  Separate upload. Not anonymized: the cover letter goes to the editor,
%  not to reviewers, so it may name the authors, the institution and the
%  public archive that the manuscript file must not name.
%
%  The journal's required cover-letter items (Cover Letter section of the
%  author instructions, retrieved 2026-08-11) are all present:
%    - title of the manuscript
%    - complete author list
%    - description of any subject overlap with previously published works
%    - description of any conflict of interest or industry support
%    - confirmation of sole submission
%  Three further items are required by the Editorial Policies page:
%    - a statement of which authors had control of the data and performed
%      the analyses (Conflicts of Interest)
%    - disclosure of prior posting, with DOI and licensing terms
%      (Preprint Servers)
%    - disclosure of how language-model tools were used, in the cover
%      letter as well as the manuscript (Large Language Models)
%
%  An earlier revision ran to four pages, most of it arguing whether a
%  Zenodo deposit counts as a preprint. That argument is not ours to
%  win, and the policy does not ask it: it asks what overlaps and how the
%  submitted manuscript differs. This version answers that with an
%  itemised overlap map and drops the rest.
%
%  One \todo{} field remains: the submission date.
%  Build: tectonic cover_letter.tex
% =====================================================================
\documentclass[11pt,letterpaper]{article}

\usepackage[T1]{fontenc}
\usepackage[utf8]{inputenc}
\usepackage{mathptmx}
\usepackage[margin=1in]{geometry}
\usepackage{microtype}
\usepackage{parskip}
\usepackage[hidelinks]{hyperref}

\pagestyle{empty}
\setlength{\emergencystretch}{2em}

\newcommand{\todo}[1]{\textbf{[TO SUPPLY: #1]}}

\begin{document}

\noindent
Sathvik Loke\\
Illinois Mathematics and Science Academy\\
1500 Sullivan Road, Aurora, IL 60506, USA\\
\texttt{sloke@imsa.edu}

\vspace{0.8em}
\noindent \todo{submission date}

\vspace{0.8em}
\noindent
The Editor\\
\emph{Radiology: Artificial Intelligence}\\
Radiological Society of North America

\vspace{0.8em}
\noindent Dear Editor,

We submit for consideration as Original Research: \textbf{``What a Slice-Level
Benchmark Certifies without the Pixels: An Audit of Six Public Imaging
Datasets,''} by Sathvik Loke, Ethan Johnson, Neeraj Movva and Aditya Raut.

The paper measures how much of a published slice-level result on a public
volumetric imaging benchmark can be reached with no images, from the subject
identifier, slice index and per-slice label alone. On the RSNA 2019 Intracranial
Hemorrhage benchmark, over 24 independently drawn patient-disjoint holdouts, one
pixel-blind score vector reads AUC 0.7381 per slice and 0.4551 per patient, while
a constant predictor reads exactly 0.500 at both units in every draw. That
comparison is between two readings of the same score vector, so it does not
depend on any published number being correct. We then identify the mechanism
rather than infer it: mean aggregation carries stack depth into the per-patient
score, and with depth held fixed the same vector reads 0.5037 across the family,
covering chance.

\subsection*{Prior public posting, and exactly what overlaps}

Before preparing this submission we deposited a research compendium on Zenodo:
concept DOI \texttt{10.5281/zenodo.21814952}, licensed CC~BY~4.0. It contains
analysis code, per-benchmark outputs, and the source of a longer manuscript.
Zenodo deposits cannot be unpublished, so we disclose it rather than offer to
remove it. If the paper is accepted we will update the record with the
publication reference and DOI, as the policy requires.

Because the policy asks how the submitted manuscript differs, here is the map
rather than a characterisation of it:

\begin{itemize}\setlength{\itemsep}{0.15em}
\item \textbf{Shared.} The benchmark set, the positional-baseline construction,
and the observation that slice-level and patient-level readings of one score
vector diverge on RSNA ICH.
\item \textbf{New here, and not in the deposit.} The estimator: the deposited
manuscript reports pooled out-of-fold cross-validation throughout, under which
the constant predictor reads 0.492 rather than 0.500---a ranking artefact of fold
identity---whereas every number in this submission is recomputed on frozen
subject-disjoint holdouts with a single fit. The stack-depth mechanism, the
depth-conditioned reading and its permutation null are new. The 24-holdout
family, which is the primary estimate here, is new. The bin-count $\times$
aggregation sensitivity analysis is new. Two arms were rebuilt from their public
source releases and re-scored; two others were withdrawn as irreproducible.
\item \textbf{In the deposit, deliberately not here.} A k-space phase study and
the other components of the larger project, one of which is being substantially
reworked. No claim in this manuscript rests on any of them.
\end{itemize}

We will send you the full deposit, or any part of it, on request.

\subsection*{The central construction is prior art, and it is not ours}

Fitting a model to a slice's position in its stack and comparing it with a
published slice-level number has been done before, and the manuscript says so in
its Introduction rather than late in the Discussion: a 2020 Kaggle notebook on
the RSNA-STR Pulmonary Embolism challenge; Yan et al's own ``Baseline: Location
feature'' row on DeepLesion; Badgeley et al on confounder-balanced test sets; and,
closest of all, Wallis and Buvat, who classify brain tumours to high accuracy
with no information about the tumour. What we add is narrower than the
construction: a within-benchmark result that needs no published comparator, an
identified mechanism rather than a gap, and a demonstration that the slice-level
reading survives the analysis choices the baseline requires while the
patient-level one does not.

\subsection*{Two things worth knowing before you read it}

\textbf{The flagship's provenance is qualified, and the manuscript says so.} The
RSNA ICH official release carries the label but neither a subject identifier nor
a slice index; both come from a public, license-permissive, pixel-free tabular
mirror. The official label file was never obtained, so the mirror could not be
checked row for row against it. Table 2 reports the internal-consistency tests,
the power of each, and what they cannot establish: they can falsify a bad mapping
and never certify a good one. The primary result compares two readings of one
score vector \emph{within} the benchmark, so it requires the mirror to be
internally consistent rather than canonical. We would rather state that plainly
than have a reviewer find it.

\textbf{One journal requirement pulls against another.} The Algorithm and Code
Transparency policy asks Materials and Methods to carry a link to the code and
the revision identifier; double-anonymized review forbids exactly that, because
our repository and archive name us. We resolved it by stating in Materials and
Methods that the tool is released with the address and revision withheld for
review and supplied on acceptance. The identifiers are above, so you have them
now; tell us if you would prefer them written into the manuscript instead.

\subsection*{Declarations}

\textbf{Sole submission.} This manuscript is submitted only to \emph{Radiology:
Artificial Intelligence}, is not under consideration at any other journal, and
will not be submitted elsewhere while it is under consideration here.

\textbf{Subject overlap.} No author has previously published on any patient,
animal or experiment reported here, and none has related work under review or in
press. The patients themselves have all been reported before, and necessarily so:
this is a retrospective secondary analysis of public label files, so every
subject was first reported by the group that released the benchmark, and many
again by the comparator papers whose published numbers appear in the supplement.
Those releases and papers are cited in Materials and Methods and Results. No new
patient data were collected, no identifiable information was used, and no pixel
intensities were read in any analysis.

\textbf{Conflicts of interest and industry support.} All four authors declare no
conflicts of interest. The study received no funding, no industry support, no
equipment, no compute donation and no in-kind support of any kind. No author has
any affiliation with, or financial interest in, any organization that produced or
maintains a benchmark audited here, or whose published results are used as
comparators. Each author will submit a completed ICMJE disclosure form.

\textbf{Control of the data and the analyses.} S.L. had full control of all data
used in the study and of the information submitted for publication, and performed
the primary analyses. S.L. wrote the primary implementation and the released
audit tool. E.T.J. wrote the second implementation reported in Table 3, working
separately and sharing no code with the first. S.L. and E.T.J. had full access to
the label files. All four authors verified the reported values against the
released outputs and take responsibility for them.

\textbf{Use of large language models.} We state this fully, because the policy
asks about study use and not only about writing, and the honest answer is broader
than the usual disclosure. Large language model assistance was used in two ways:
in drafting and editing this manuscript, and in writing the analysis code that
produces every number in the Results. The tool was \textbf{Claude (Opus~5),
Anthropic}, accessed 2026-07-27 to 2026-08-14.

No value reported in the paper was generated by a language model. Every number is
a deterministic output of the released code, run on a published label file, and
the flagship estimate was reproduced by a second implementation written
separately from the primary one and sharing no code with it, agreeing to 0.0003
AUC per slice and 0.005 per patient with a constant predictor of exactly 0.500 in
every draw of both. The two implementations were written by different authors,
S.L. and E.T.J., so the check can catch an implementation error in either. We
still describe it as a code-independence check rather than an independent
replication: both authors used the assistance disclosed above, so it cannot catch
a misconception about the data that both shared. That check, rather than any
assurance about provenance, is what should persuade a reviewer, and the code is
released so it can be repeated without us.

The authors have read, verified against the released artefacts, and take full
responsibility for every claim in the manuscript.

\textbf{Checklist.} A CLAIM checklist is uploaded with this submission. The study
audits an evaluation protocol rather than developing a diagnostic model, so a
number of CLAIM items concern steps that were not performed and are marked not
applicable, each with a reason. If you would prefer STROBE for a study of this
shape, we will complete and send it.

\subsection*{Reviewers}

Reviewers with expertise in benchmark and challenge design, in meta-research and
reporting standards, or in neuroradiology and head CT triage would be well placed
to assess this work. We request no exclusions.

Thank you for considering the manuscript.

\vspace{1.0em}
\noindent Yours sincerely,

\vspace{2.0em}
\noindent
Sathvik Loke\\
on behalf of Sathvik Loke, Ethan Johnson, Neeraj Movva and Aditya Raut

\end{document}
